\section{Verification purpose}

\textbf{Case ID:} \texttt{\CaseID}. This case exercises simultaneous spatial fuel variation, opposing slopes, and steady wind to reveal integration failures that isolated cases may miss.  It is currently diagnostic rather than decisional because no independent quantitative reference or tolerance has been established.

\section{Mathematical formulation and numerical implementation}

At each active front cell ELMFIRE combines fuel-dependent reaction terms, wind and slope factors, and elliptical propagation geometry to obtain a local normal velocity.  The configuration intentionally produces competing directional effects.  No closed-form solution exists for the discontinuous, coupled field.

\section{Derivation of expected behavior}

The expected behavior is an invariant ordering rather than a numerical solution: spread should remain finite and physically coherent, respond to the four fuel regions, accelerate on the steeper slope, and show net wind alignment.  Because those statements have not yet been translated into independent numerical bounds, they cannot support PASS.

\section{Verification metrics and acceptance criteria}

The current postprocessor records maximum finite spread rate as a diagnostic.  It is useful for detecting missing or explosive output but is insufficient to verify spatial interaction.  The case is therefore reported as NOT EVALUABLE even when output exists.  A future contract should add regional arrival-time or directional-front metrics with justified bounds.

\section{Simulation configuration and predicted outcome}

The 30 m domain combines the CASE23 fuel quadrants, CASE22 opposing slopes, 5 mph wind, uniform moisture, and central ignition. Predicted output is asymmetric, spatially continuous, and preferentially extended by both wind and locally faster fuels/slopes.

\section{Actual simulation results}

The required evidence is generated only from case-local ELMFIRE outputs. The current postprocessing status is \textbf{\VerificationStatus}; \PassingMetricCount{} of \MetricCount{} required metric rows pass across \VariantCount{} variant(s). A missing output is reported as NOT RUN or NOT EVALUABLE and is not interpreted as verification evidence.


\begin{figure}[H]
  \centering
  \IfFileExists{../figures/whole_domain_result.pdf}{%
    \includegraphics[width=0.82\textwidth,height=0.58\textheight,keepaspectratio]{whole_domain_result.pdf}%
  }{\fbox{\parbox{0.85\textwidth}{No case-local whole-domain ELMFIRE raster is available.}}}
  \caption{Representative whole-domain case-local ELMFIRE output. The generated figure records its source path, coordinate reference system, and variant.}
\end{figure}


\begin{figure}[H]
  \centering
  \IfFileExists{../figures/verification_summary.pdf}{%
    \includegraphics[width=\textwidth,height=0.58\textheight,keepaspectratio]{verification_summary.pdf}%
  }{\fbox{\parbox{0.85\textwidth}{No case-local summary figure is available.}}}
  \caption{Case-local expected and measured verification metrics. A missing measured value denotes unavailable or incomplete ELMFIRE output.}
\end{figure}

\section{Calculated metrics and verification decision}

\begin{center}
\begin{tabular}{@{}llll@{}}
\toprule
Variant / metric & Expected & Calculated & Status \\
\midrule
\MetricRows
\bottomrule
\end{tabular}
\end{center}

The authoritative overall decision is \textbf{\VerificationStatus} and is generated by \texttt{scripts/postprocess.py} from \texttt{outputs/metrics.json}; it is not recomputed in this report.

\section{Technical interpretation}

Plausible appearance is not a verification decision. Any interpretation remains diagnostic until independent quantitative metrics are defined. No scientific conclusion should be drawn when the status is NOT RUN or NOT EVALUABLE.

\section{Scope and limitations}

No quantitative acceptance rule currently exists. This case must not be counted as a passing verification case. This case provides implementation-verification evidence only.
